\documentclass[12pt]{article}
\usepackage{physics}

\begin{document}
Over last week's lab session and my work between the first and second lab sessions, I was able to design and experimentally test/validate the Tri-Channel Automatic Thermometer (TCAT), an automatic, parallelized data collection tool for measuring the cooling curves of all three pipes simultaneously so that we can save time during the lab. I was able to validate that this device produces accurate readings, that it is consistent and reliable, and that it is an overall effective method of data collection for this lab. Then, using the TCAT, Eric and I were able to take a first round of ``official'' data (meaning data where the emissivities produced from fitting can be reliably trusted and added to a final backlog of emissivities used to produce the final approximations), but due to our differential equation solvers being incomplete at the time of the lab, we only intended to complete a qualitative analysis of the emissivities by using the TCAT's parallelization to compare their cooling curves, and we determined that their characteristic cooling curves match literature predictions (i.e., the emissivity of the lacquered pipe must be the greatest and the emissivity of the smooth pipe must be the lowest).

Moving into next week, the original plan was to perform our first quantitative analysis of the data we've collected (i.e., performing a fit and obtaining emissivities), but shortly after the lab concluded, I realized why my differential equation solver was not working (I had been using the Boltzmann constant $k_B$ rather than the Stefan-Boltzmann constant $\sigma$ on accident), and I began fitting right away. Therefore, my main goal is to perform a second round of data collection and refit the data to compare the emissivities and compare the results between the first and second rounds. The overarching goal moving forward is to perfect the testing process, so I'd like to also try seeing if dynamic time spacing could be built into the TCAT to dynamically collect more data at higher temperatures and less data at lower temperatures (as the current data volume is higher than necessary, and makes communicating uncertainty difficult).

Lastly, once again, our quantity of interest is emissivity $\varepsilon$. We're using a non-analytical, numerical fit based on building a differential equation and solving it via a forward model. The model is essentially based on the radiation-convection heat transfer model
\begin{equation}
  \dv{T}{t} = \frac{1}{\rho V C} \left[ \dv{Q}{t}_{rad} + \dv{Q}{t}_{conv} \right] 
\end{equation}
which is numerically solved with \texttt{solve\char`_ivp} given the initial condition $T(t_0) = T_0$ to produce the numerically-solved general solution $T_\varepsilon(t)$ (thus the $x$-axis would be time and the $y$-axis would be tmperature), and this can be varied with \texttt{curve\char`_fit} to calculate $\varepsilon \pm \delta\varepsilon$ as $\dv{Q}{t}_{conv}$ is a function of $\varepsilon$. With this, we don't intend to use $x$-uncertainties (due to manually-setting the timescale on the TCAT, so this uncertainty is instrumental, uniform across all measurements, and essentially negligible), and the $y$-uncertainties are instrumental and uniform across all measurements due to being the result of the imprecision of the MAX6675 thermocouple amplification chip on which the TCAT is comprised. 
\end{document}
